%
\vspace*{0.1cm}
\begin{center}
{\bf\Large epic-[AN/TN]-[TC/SC/AC]-[Year]-[\#] } \\
\vspace*{1.0cm}
{\bf\Large LightGBM-Guided Exploration of a Six-Dimensional ePIC nHCal Design Space for Low-Energy Neutron Detection} \\
\end{center}
%
\vspace*{0.1cm}
%
\begin{center}
{\bf Principal Author List:} \\
\vspace*{0.15cm}
Author 1 (Institution, Email)\\
Author 2 (Institution, Email)\\
Author 3 (Institution, Email)\\
\vspace*{0.5cm}
(\today)
\end{center}



\vfill
\begin{center}
{\includegraphics[width=0.4\linewidth]{2560px-EPIC-logo_black_transparent.png}}
\end{center}
\vfill


\pagebreak
%
\begin{abstract}
\noindent \textit{The ePIC negative-eta hadronic calorimeter (nHCal) is a backward sampling calorimeter intended to improve hadronic coverage in the electron-going region, including neutral-hadron detection and hit-level activity measurements. This note tests whether the simulated low-energy neutron hit-detection efficiency can be improved by varying the absorber and scintillator thicknesses in three longitudinal nHCal segments, forming a six-dimensional geometry design space. Candidate geometries were generated with a standalone DD4hep/DDsim workflow and simulated with Geant4 using a neutron G4GPS source spanning 0.1--2.5 GeV/c. A LightGBM surrogate-guided search was used to rank geometries using neutron detection efficiency, average fired-tile count, and average fired-layer count. The scanned landscape was shallow: the final selected geometry improved the baseline efficiency from $0.53172 \pm 0.00331$ to $0.54184 \pm 0.00276$, while fired-tile and fired-layer activity increased by about 4\%. These results demonstrate that surrogate-guided optimization can be used to explore the nHCal geometry design space and identify geometry-dependent changes in neutron hit-detection efficiency. Within the tested bounds, however, the best confirmed gain is modest. This suggests that the current direction can identify small gains, but higher-statistics scans and additional systematic studies would be needed to determine whether stronger improvements exist elsewhere in the design space.
}
\end{abstract}
